\chapter*{INTRODUCCIÓN}
\addcontentsline{toc}{chapter}{INTRODUCCIÓN}
\thispagestyle{empty}

En el Perú, entre 2010 y 2025, la máxima demanda del Sistema Eléctrico Interconectado Nacional (SEIN) aumentó de \SI{4579}{\mega\watt} a \SI{7942}{\mega\watt} y la producción eléctrica nacional alcanzó \SI{65109}{\giga\watt\hour} al cierre preliminar de 2025 \parencites[sec.~10.4]{minem2025anuario2024}[cuadro~4]{minem2025indicadoresDiciembre}{coesMaximaDemanda}. En ese mismo período, la participación conjunta de la generación solar y eólica pasó desde un valor cercano a cero hasta 9,5\,\pct{}. La incorporación de generación distribuida geográficamente exige transportar energía entre nuevos puntos y aprovechar la capacidad disponible de la red \parencite{iea2025transmissiongrid}.

Los cables enterrados se utilizan en redes colectoras internas de plantas solares y eólicas, así como en enlaces urbanos donde el espacio, la servidumbre, la seguridad o la integración arquitectónica dificultan construir líneas aéreas \parencites{cigre2023overheadLines}{coes2021puntaLomitasEO}{coes2022planTransmision2023_2032}{proinversion2026palian}. La capacidad de transporte de estas líneas depende, además de las características propias del conductor y su aislamiento, de la evacuación del calor generado hacia el lecho de asiento, el suelo de relleno y el suelo nativo. La ampacidad $\Imax$ es la corriente que un cable puede transportar de manera sostenida, bajo las condiciones de instalación especificadas, sin que la temperatura máxima del conductor $\Tmax$ supere el límite térmico admisible $T_{\mathrm{lim}}$ \parencites[cláusulas~4--5]{iec60287}[3--5]{enescu2020}.

El objeto de estudio de la investigación es el arreglo de cables--entorno térmico, integrado por los cables multicapa, su disposición espacial y los materiales circundantes que determinan la transferencia de calor \parencites[1--3]{enescu2020}[41--42]{cigre2025}. Este entorno resulta heterogéneo cuando su conductividad térmica $\kx$ varía espacialmente por humedad, estratos, rellenos térmicos o interfaces entre materiales. Dicha variación altera el campo de temperatura $T(x,y)$, a partir del cual se determinan $\Tmax$ e $\Imax$ \parencites[3--4]{enescu2021}[1,6--7]{khumalo2025}[3]{xing2023}.

El problema técnico aparece cuando una propiedad homogénea equivalente $k_{eq}$ sustituye esa distribución espacial en el cálculo térmico. Esta simplificación puede ocultar zonas de baja disipación y subestimar la temperatura real, o bien conducir a estimaciones conservadoras que restringen la capacidad operativa. Los métodos numéricos de campo, como el método de elementos finitos (FEM), permiten representar la geometría y los materiales con detalle, pero exigen construir, mallar y verificar cada escenario por separado \parencites[41--42,93--99]{cigre2025}[1--3,12]{aldulaimi2024}.

La investigación desarrolló y evaluó un artefacto basado en una red neuronal informada por física (PINN) que resuelve la ecuación de conducción de calor estacionaria bidimensional y sus condiciones de frontera e interfaz \parencite{raissi2019}. El artefacto aproxima $T(x,y)$ en el conductor, las capas concéntricas y el suelo, y extrae la temperatura máxima del dominio interior sin emplear soluciones FEM como etiquetas de entrenamiento. El cálculo de $\Imax$ requiere resolver adicionalmente la condición de corriente límite con la misma ley de pérdidas eléctricas. Se evaluaron distintas representaciones, arquitecturas por subdominios, presupuestos y distribuciones de puntos de colocación, constatando que una red mayor o una pérdida de entrenamiento reducida no bastan para acreditar exactitud frente a una referencia convergente. La ciencia del diseño (DSR) estructuró la relación entre el problema práctico, la construcción iterativa del artefacto y la evidencia empírica obtenida \parencites[54--56]{peffers2007}[342--350]{gregor2013}.

El documento sigue la estructura prescrita para tesis de maestría profesionalizante o de especialización de la Universidad Nacional de Ingeniería. El Capítulo~\ref{ch:planteamiento} presenta el diagnóstico, identifica y formula el problema, establece los objetivos y las proposiciones de diseño, define la operacionalización discreta de variables y detalla la metodología DSR. El Capítulo~\ref{ch:marco-teorico} reúne los antecedentes bibliográficos en orden cronológico y temático, desarrolla las bases físicas, numéricas y metodológicas del problema y precisa el glosario de términos técnicos. El Capítulo~\ref{ch:propuesta-desarrollo} expone la propuesta computacional, el desarrollo secuencial de las etapas A a D, el análisis de resultados, la discusión e interpretación de los hallazgos y la estimación del impacto técnico, de riesgos y de sostenibilidad. El Capítulo~\ref{ch:conclusiones} establece las conclusiones y recomendaciones correspondientes a cada objetivo específico. Finalmente, las referencias documentan las fuentes citadas y los anexos presentan la matriz de consistencia, el protocolo preliminar de juicio de expertos, las matrices de operacionalización y trazabilidad, el catálogo de casos, el protocolo de reproducción y el expediente completo de ejecuciones.
